% Audited research fragment: angular prime separation under Dirichlet GRH.
% Written 2026-11-26. This is a fragment, not a standalone LaTeX document.

\section{A uniform angular prime-separation lemma}

All logarithms below are natural.  Zeros are always counted with their
multiplicities and with both signs of their ordinates.  We assume Dirichlet
GRH throughout this section, including RH for the Riemann zeta function.

\begin{lemma}[A coarse zero count proved directly from GRH]
Let $\chi$ be a primitive nonprincipal character of conductor $q>1$, and let
$N_\chi(T)$ count its nontrivial zeros with $|\gamma|\le T$.  If $T\ge2$, then
\[
 N_\chi(T)\le T\log(qT).
\]
For the Riemann zeta function,
\[
 N_\zeta(T)\le2T\log(2T)\qquad(T\ge2).
\]
Consequently the common, deliberately loose envelope
\begin{equation}\label{eq:angular-count}
 N_\chi(T)\le2T\log(2qT)\qquad(T\ge2)
\end{equation}
holds for primitive characters, with $q=1$ representing zeta.
\end{lemma}

\begin{proof}
Put $a=(1-\chi(-1))/2\in\{0,1\}$ and
\[
 \Lambda_\chi(s)=(q/\pi)^{(s+a)/2}
                    \Gamma((s+a)/2)L(s,\chi).
\]
For a nonprincipal primitive character this is entire of order one.  Its
functional equation gives
$|\Lambda_\chi(1/2+u)|=|\Lambda_\chi(1/2-u)|$ for real $u$.
The centered Hadamard product, together with GRH, therefore gives
\begin{equation}\label{eq:angular-logder}
 \Re\frac{\Lambda_\chi'}{\Lambda_\chi}(1/2+u)
       =\sum_\gamma\frac{u}{u^2+\gamma^2}\qquad(u>0).
\end{equation}
For completeness, write the centered product as
$z^r e^{A+Bz}\prod_{\gamma\ne0}(1-z/(i\gamma))e^{z/(i\gamma)}$.
On the real axis the product factors have modulus
$(1+z^2/\gamma^2)^{1/2}$, while the auxiliary exponentials have modulus one.
The displayed modulus symmetry forces $\Re B=0$; differentiating its
logarithm proves \eqref{eq:angular-logder}, including the term $r/u$ for
central zeros.  The series converges absolutely.

Every zero with $|\gamma|\le T$ contributes at least $1/(2T)$ at $u=T$.
For $\sigma=T+1/2\ge5/2$, the Euler product and an elementary integral
comparison give
\[
 \left|\frac{L'}{L}(\sigma,\chi)\right|
 \le\sum_{n=2}^\infty\frac{\log n}{n^{5/2}}
 \le\frac1{\sqrt2}\left(\frac7{12}\log2+\frac29\right)<\frac12.
\]
The inequality $\psi_\Gamma(y)\le\log y$ for $y>0$ follows, for example,
from Jensen's inequality applied to the logarithm of a gamma random variable
of mean $y$.  Thus
\begin{align*}
 \Re\frac{\Lambda_\chi'}{\Lambda_\chi}(T+1/2)
 &\le\frac12\log\frac q\pi
       +\frac12\log\frac{T+a+1/2}{2}+\frac12\\
 &\le\frac12\log(qT)
       +\frac12\log\frac{7e}{8\pi}
 <\frac12\log(qT).
\end{align*}
Here $T\ge2$, $a\le1$, and $7e<8\pi$.  This proves the first bound.

For zeta use the entire function
$\xi(s)=\tfrac12s(s-1)\pi^{-s/2}\Gamma(s/2)\zeta(s)$ in the same
centered-product argument.  Since $\zeta'/\zeta(\sigma)<0$ for $\sigma>1$,
\begin{align*}
 \frac{\xi'}{\xi}(T+1/2)
 &\le\frac1{T+1/2}+\frac1{T-1/2}
          +\frac12\log\frac{T+1/2}{2\pi}\\
 &\le\frac{16}{15}+\frac12\log T
          +\frac12\log\frac5{8\pi}
 \le\log(2T).
\end{align*}
For the final inequality, $\log(5/(8\pi))<-1$,
$17/30<\log2$, and $T\ge1$ suffice.  Again
$N_\zeta(T)\le2T\xi'/\xi(T+1/2)$.
\end{proof}

\begin{lemma}[An order-independent angular prime]
Let $M\ge3$ and let $\psi$ be a nonprincipal Dirichlet character modulo $M$,
not necessarily primitive.  Set
\[
 L=\log(2M),\qquad x=10^6 L^2.
\]
There exists a prime $p$ such that
\begin{equation}\label{eq:angular-prime}
 \frac{x}{e}\le p\le ex<3\cdot10^6L^2,
 \qquad p\nmid M,\qquad \Re\psi(p)\le\frac12.
\end{equation}
In particular $|1-\psi(p)|\ge1$.  The constants do not depend on the
order of the character.
\end{lemma}

\begin{proof}
Use the triangular function and its Fourier transform
\[
 g(v)=(1-|v|)_+,\qquad
 h_x(t)=e^{it\log x}\left(\frac{\sin(t/2)}{t/2}\right)^2,
\]
where the quotient at zero is defined as $1$.
Write
\[
 W_\chi(x)=\sum_{n\ge2}\frac{\Lambda(n)\chi(n)}{\sqrt n}
                   g(\log(n/x)).
\]
For a primitive character of parity $a$, Weil's explicit formula, with
the Fourier convention $h(t)=\int_{\mathbb R}f(v)e^{itv}\,dv$, gives
\begin{equation}\label{eq:angular-explicit}
 W_\chi(x)=\mathbf1_{q=1}\,c(\sqrt x+x^{-1/2})
               -G_a(x)-\sum_\gamma h_x(\gamma),
 \quad
 c=\left(\frac{\sinh(1/4)}{1/4}\right)^2>1,
\end{equation}
where
\[
 G_a(x)=\int_{\log x-1}^{\log x+1}
    g(v-\log x)\frac{e^{-(a+1/2)v}}{1-e^{-2v}}\,dv.
\]
There is no conductor term because the test function in the $v$ variable
vanishes at zero, and no negative-frequency prime term because its support
is positive.  Formula \eqref{eq:angular-explicit} can equally be obtained
first for smooth approximations to $g$ and then by dominated convergence;
the transform decays as $O((1+|t|)^{-2})$ and the zero sum is absolutely
convergent.  Since $x\ge e^2$,
\[
 0\le G_a(x)\le\frac{e^{1/2}}{\sqrt x(1-e^2/x^2)}
                         <\frac2{\sqrt x}<1.
\]

The zero count \eqref{eq:angular-count} gives an explicit uniform estimate.
The zeros with $|\gamma|\le2$ contribute at most $4\log(4q)$.
For the others, partial summation gives
\begin{align*}
 \sum_{|\gamma|>2}|h_x(\gamma)|
 &\le4\sum_{|\gamma|>2}\gamma^{-2}
 \le8\int_2^\infty N_\chi(t)t^{-3}\,dt\\
 &\le16\int_2^\infty\frac{\log(2qt)}{t^2}\,dt
 =8\bigl(\log(4q)+1\bigr).
\end{align*}
Thus we may use the convenient bound
\[
 B(q)=20\log(2q)+21
\]
to conclude
\[
 |W_\chi(x)|\le B(q)\quad(q>1),\qquad
 W_{\mathbf1}(x)\ge c\sqrt x-B(1)\quad(q=1).
\]

Let $\psi^*$ be the primitive character inducing $\psi$, of conductor
$d\mid M$.  Since $\psi$ is nonprincipal, so is $\psi^*$.
Removing Euler factors at primes dividing $M$ changes the weighted sum by
at most
\begin{align*}
 E_{\rm ram}
 &\le\sum_{p\mid M}\log p\sum_{k\ge1}p^{-k/2}
 =\sum_{p\mid M}\frac{\log p}{\sqrt p-1}
 <\frac52\log M.
\end{align*}
Consequently
\[
 |W_\psi(x)|\le B(d)+E_{\rm ram},\qquad
 W_{\chi_{0,M}}(x)\ge c\sqrt x-B(1)-E_{\rm ram},
\]
where $\chi_{0,M}$ is the principal character modulo $M$.

The total contribution of all prime powers of exponent at least two is
bounded by
\begin{align*}
 E_{\rm pp}
 &\le\sum_{p\le\sqrt{ex}}\log p\sum_{k\ge2}p^{-k/2}
 \le4\sum_{2\le n\le\sqrt{ex}}\frac{\log n}{n}\\
 &\le\frac12(1+\log x)^2+4.
\end{align*}
Here
$\sum_{2\le n\le y}\log n/n\le\tfrac12\log^2y+1$
follows by separating the first two terms and comparing the rest with the
integral of $\log t/t$; alternatively its unimodality gives the same
integral bound with a maximum-term error smaller than $1$.

Suppose that every prime in the support of $g(\log(p/x))$, coprime to $M$,
has $\Re\psi(p)>1/2$.  Positivity of the weights and the estimate
$\Re\psi(p^k)\ge-1$ imply
\[
 \Re W_\psi(x)\ge\frac12W_{\chi_{0,M}}(x)
                                      -\frac32E_{\rm pp}.
\]
Combining the last three bounds yields
\begin{align*}
 c\sqrt x
 &\le2B(d)+B(1)+3E_{\rm ram}+3E_{\rm pp}\\
 &\le47.5L+89+1.5(1+\log x)^2.
\end{align*}
This contradicts $c\sqrt x>1000L$.  Indeed, with
$A=1+\log10^6<15$, the function
$(A+2\log L)^2/L$ decreases for $L\ge\log6$, because its derivative has
the sign of $4-(A+2\log L)$.  Its maximum on this range is less than
$16.2^2/1.7<155$.  Also $89/L<53$.  The last upper bound is therefore
less than $(47.5+53+1.5\cdot155)L=333L<1000L$.
The contradiction proves \eqref{eq:angular-prime}.  Finally, for
$|\psi(p)|=1$, $|1-\psi(p)|^2=2-2\Re\psi(p)\ge1$.
\end{proof}

\begin{corollary}[Separating arbitrary complex characters]
Let $\chi_1,\chi_2$ be distinct primitive characters of conductors
$q_1,q_2\le Q$, and let $M=\operatorname{lcm}(q_1,q_2)\ge3$.
Then there is a prime $p\nmid M$ in the annulus from
\eqref{eq:angular-prime} such that
\[
 |\chi_1(p)-\chi_2(p)|\ge1,
 \qquad
 p<12\cdot10^6\log^2(2Q).
\]
\end{corollary}

\begin{proof}
The quotient $\psi=\chi_1\overline{\chi_2}$, viewed modulo $M$, is
nonprincipal: equality of the induced characters would force equality of
their unique primitive inducing characters.  Apply the lemma and use
$\log(2M)\le2\log(2Q)$, since $M\le Q^2$.
\end{proof}
